% TOP_PROBLEM: {"id": 7200010, "title": "In spherical or hyperbolic geometry, must polyhedra with the same volume and Dehn invariant be scissors-congruent", "category_id": 5, "category": {"id": 5, "name": "algebraic_geometry", "display_name": "Algebraic Geometry", "description": "Geometric objects defined by polynomial equations.", "slug": "algebraic-geometry", "order_index": 5, "created_at": "2026-07-31T15:26:25.671Z"}, "status": "partially_solved"}
% FIRST_POSTED: 2026-09-25
% ENTRY_KIND: statement_only
% !TeX program = lualatex
% Definition and sources only; this file is not an LLM solution attempt.
\documentclass[11pt]{article}
\usepackage[margin=25mm]{geometry}
\usepackage{fontspec}
\setmainfont{Latin Modern Roman}
\usepackage{amsmath,amssymb,mathtools}
\usepackage{unicode-math}
\setmathfont{Latin Modern Math}
\usepackage{xurl,hyperref}
\hypersetup{colorlinks=true,urlcolor=blue}
\setcounter{secnumdepth}{0}
\setlength{\parindent}{0pt}
\setlength{\parskip}{5pt}
\setlength{\emergencystretch}{3em}
\begin{document}
\section*{\texorpdfstring{In spherical or hyperbolic geometry, must polyhedra with the same volume and Dehn invariant be scissors-congruent}{In spherical or hyperbolic geometry, must polyhedra with the same volume and Dehn invariant be scissors-congruent}}\label{7200010}
\subsection{Definitions and mathematical statement}
Two three-dimensional polyhedra in $X=\mathbb H^3$ or $S^3$ are scissors congruent if finite decompositions into polyhedral pieces can be matched by isometries. Their classes form the scissors group $\mathcal P(X)$, imposing additivity under cuts. The Dehn invariant is the additive invariant represented, in the source's angular convention, by
\[
 D(P)=\sum_{e}\ell(e)\otimes(\theta_e\bmod\pi{\mathbb{Q}}).
\]
Here $\ell(e)$ and $\theta_e$ are edge lengths and dihedral angles; the precise target group and spherical relations are those of Neumann's formulation. The task, separately for both geometries, is injectivity of volume on $\ker D$. Equivalently, equal volume and equal Dehn invariant should imply scissors congruence. In the spherical scissors group volume uses its stated full-sphere convention; one must not replace the group or allow ideal pieces without checking those conventions.
\par\smallskip\textit{Formulation and concept references:} \href{https://arxiv.org/pdf/math/9712226}{[S1]}.

\subsection{Short English statement}
In spherical and hyperbolic three-dimensional geometry, do volume and the Dehn invariant together completely determine scissors congruence?
\subsection{Sources}
\begin{itemize}
\item[S1]Neumann, Hilbert's 3rd Problem and invariants of 3-manifolds, Geom. Topol. Monogr. 1 (1998), 383-411.
\url{https://arxiv.org/pdf/math/9712226}.
\textit{Formulation source.}
\end{itemize}
\end{document}
